\indent ({\kbvlatin i}) ನಾನು, ನಾವು—ಇವು “ಉತ್ತಮ ಪುರುಷ” ಸರ್ವನಾಮಗಳು. ({\kbvlatin ii}) ನೀನು, ನೀವು—ಇವು “ಮಧ್ಯಮ ಪುರುಷ” ಸರ್ವನಾಮಗಳು. ({\kbvlatin iii}) ಅವನು, ಅವಳು, ಅವರು, ಇವನು, ಇವಳು, ಇವರು, ಅದು, ಅವು, ಇದು, ಇವು—ಇವು “ಪ್ರಥಮ ಪುರುಷ” ಸರ್ವನಾಮಗಳು.\par

\vspace{4pt plus 4pt minus 1pt}\noindent{\kbvsecfont\fontsize{11pt}{14.5pt}\selectfont\color{kbvaccentdark} (1)\hspace{0.45em}ಆತ್ಮಾರ್ಥಕ ಸರ್ವನಾಮಗಳು :}\hspace{0.5em}ಉದಾಹರಣೆಗೆ ತಾನು, ತಾವು ಇತ್ಯಾದಿ ಆತ್ಮಾರ್ಥಕ ಸರ್ವನಾಮಗಳು.\par

\vspace{4pt plus 4pt minus 1pt}\noindent{\kbvsecfont\fontsize{11pt}{14.5pt}\selectfont\color{kbvaccentdark} (2)\hspace{0.45em}ಪ್ರಶ್ನಾರ್ಥಕ ಸರ್ವನಾಮಗಳು :}\hspace{0.5em}ಉದಾಹರಣೆಗೆ ಯಾರು, ಎಲ್ಲಿ, ಯಾವುದು, ಏನು ಇತ್ಯಾದಿ ಪ್ರಶ್ನಾರ್ಥಕ ಸರ್ವನಾಮಗಳು.\par

\kbvexercise{ಅಭ್ಯಾಸ ಪತ್ರಿಕೆ ೮}

\vspace{2.5pt plus 3pt minus 0.5pt}\noindent\hspace*{1.2em}\kbvnum{(1)}\hspace{0.55em}ವಸ್ತುವಾಚಕಗಳು ಎಂದರೇನು ? ಉದಾಹರಣೆಗಳ ಸಹಿತ ತಿಳಿಸಿ.\par

\vspace{2.5pt plus 3pt minus 0.5pt}\noindent\hspace*{1.2em}\kbvnum{(2)}\hspace{0.55em}ವಸ್ತುವಾಚಕದ ಮೂರು ಪ್ರಕಾರಗಳನ್ನು ತಿಳಿಸಿ.\par

\vspace{2.5pt plus 3pt minus 0.5pt}\noindent\hspace*{1.2em}\kbvnum{(3)}\hspace{0.55em}ರೂಢನಾಮಗಳಿಗೂ ಅಂಕಿತನಾಮಗಳಿಗೂ ಇರುವ ವ್ಯತ್ಯಾಸವೇನು ? ಉದಾಹರಣೆಗಳ ಸಹಿತ ತಿಳಿಸಿ.\par

\vspace{2.5pt plus 3pt minus 0.5pt}\noindent\hspace*{1.2em}\kbvnum{(4)}\hspace{0.55em}ಗುಣವಾಚಕ ಎಂದರೇನು ? ಉದಾಹರಣೆಗಳ ಸಹಿತ ವಿವರಿಸಿ.\par

\vspace{2.5pt plus 3pt minus 0.5pt}\noindent\hspace*{1.2em}\kbvnum{(5)}\hspace{0.55em}ಸುಯೋಧನ, ದಶರಥ, ಸುಬಾಹು, ದ್ರೌಪದಿ—ಈ ಪುರಾಣ ಪ್ರಸಿದ್ಧ ಹೆಸರುಗಳು ಯಾವ ಪ್ರಕಾರದ ನಾಮವಾಚಕಗಳು?\par

\vspace{2.5pt plus 3pt minus 0.5pt}\noindent\hspace*{1.2em}\kbvnum{(6)}\hspace{0.55em}ಸಂಖ್ಯಾವಾಚಕಗಳಿಗೂ ಸಂಖ್ಯೇಯವಾಚಕಗಳಿಗೂ ಇರುವ ವ್ಯತ್ಯಾಸವನ್ನು ಉದಾಹರಣೆಗಳ ಮೂಲಕ ಸ್ಪಷ್ಟಪಡಿಸಿ.\par

\vspace{2.5pt plus 3pt minus 0.5pt}\noindent\hspace*{1.2em}\kbvnum{(7)}\hspace{0.55em}ಈ ಪದಗಳನ್ನು ಬಳಸಿಕೊಂಡು ಸ್ವಂತ ವಾಕ್ಯಗಳನ್ನು ರಚಿಸಿ : ಅಷ್ಟೊಂದು, ಇಂತಿಷ್ಟು, ಇಂಥಾದ್ದು, ಎಂತಹ\par

\vspace{2.5pt plus 3pt minus 0.5pt}\noindent\hspace*{1.2em}\kbvnum{(8)}\hspace{0.55em}ಇವುಗಳಿಗೆ ಎರಡೆರಡು ಉದಾಹರಣೆಗಳನ್ನು ಕೊಡಿ : ರೂಢನಾಮಗಳು, ಪರಿಮಾಣವಾಚಕಗಳು ಅಂಕಿತನಾಮಗಳು, ಪ್ರಕಾರವಾಚಕಗಳು ಅನ್ವರ್ಥನಾಮಗಳು, ಪ್ರಥಮ ಪುರುಷ ಸರ್ವನಾಮಗಳು\par
